% figuras/tikz/cap03-omnicanal.tex
% Experiencias omnicanal: puntos de contacto distintos que deben resolver a
% una sola identidad de cliente (Verhoef, Kannan & Inman, 2015).
\begin{tikzpicture}[
    canal/.style={rectangle, rounded corners, draw=bronceIPN, thick, fill=bronceIPNClaro,
                  minimum width=2.4cm, minimum height=0.9cm, align=center, font=\scriptsize},
    centro/.style={circle, draw=guindaIPN, very thick, fill=guindaIPNClara,
                   minimum size=2.4cm, align=center, font=\small},
    flecha/.style={-{Stealth[length=2.5mm]}, thick, draw=doradoIPN}
  ]
  \node[centro] (id) at (0,0) {Identidad\\ única de\\ cliente};

  \node[canal] (app)    at (-4.2, 2.2) {App móvil};
  \node[canal] (web)    at (0, 3.2)    {Sitio web};
  \node[canal] (tienda) at (4.2, 2.2)  {Tienda física};
  \node[canal] (social) at (-4.2,-2.2) {Redes sociales};
  \node[canal] (cc)     at (4.2,-2.2)  {Centro de contacto};

  \foreach \n in {app,web,tienda,social,cc} {
    \draw[flecha] (\n) -- (id);
  }
\end{tikzpicture}
